\Qhead{5.2}{Распространение радиоволн}{Геометрическое и дифракционное приближения при анализе распространения радиоволн. Влияние неровностей земной поверхности. Земные и тропосферные радиоволны. Рассеяние и поглощение радиоволн в тропосфере. Эффект "замирания". Тропосферный волновод. Распространение радиоволн в ионосфере. Дисперсия и поглощение радиоволн в ионосферной плазме. Ионосферная рефракция. Ход лучей в подводном звуковом канале и тропосферном радиоволноводе.}
\Subfile{5.2.1}{sections/05_antennas/q02/s01}
\Subfile{5.2.2}{sections/05_antennas/q02/s02}
\Subfile{5.2.3}{sections/05_antennas/q02/s03}
\Subfile{5.2.4}{sections/05_antennas/q02/s04}

\recap{Лучи описывают плавную рефракцию и многолучевость; дифракция заполняет геометрическую тень. Поверхность, тропосфера и плазма создают разные потери и каналы. Звуковой и радиоканал объединяет наличие поворотных лучей.}
\board{Первая зона Френеля; сумма многолучевых полей; профиль $n$ или $c$; условие поворота и ионосферная проницаемость.}
\oral{Почему ионосферный луч возвращается на Землю без отражения от резкой металлической границы?}
\bait{Можно оставить различие фазового и группового распространения в плазме: их скорости различны, но не нарушают причинность.}
\sources{Сорокин~\cite{sorokinray,sorokinsignal}, теоретические разделы; Цысарь~\cite{tsysar1}, слоистая среда. Плазменная модель сверена по Fitzpatrick~\cite{fitzpatrick}; звуковой канал выведен из сохранения горизонтальной медленности.}
